% SPDX-License-Identifier: Apache-2.0
\section{An Output Held Across States}
\label{sec:x-hold}

The run's wire keeps its value between sets. An output set in one
state and left alone in the next is still what it was, and a
protocol leans on that: a \code{busy} raised when work is taken and
dropped when it is done is set in two states and read in all of
them. A wire of the netlist has no memory, so a process of several
waits cannot drive such an output as a wire selected on the state.

The lowering holds it instead. Each output a state sets gets a
register the unit does not declare, \code{<port>\_held}, driven
inside the state's arm of the clocked block, so it takes the state's
value at the edge that leaves the state and keeps it until another
state's edge. The port itself is the state's expression while the
state is about to leave, that is while the register holds the state's
number and the wait's condition holds, and the register otherwise.
The two halves are one wire, \code{cond ? value : held}, chained
over the states that set the port.

Why both halves, when the register alone would hold the value: the
netlist's wires run a cycle ahead of the run's. The run sets a wire
at the edge from what stood before it, and the netlist's wire is
computed from what stands now, which is the same values one cycle
earlier; the testbench compares them with that offset. A register
alone would be a cycle late at every change, and a wire alone would
show the next state's value while the wait stalled. The combination
is exact, and the build checks it: the stepper below is
\code{ex\_seq}'s with a \code{phase} output set in two of its three
states and a \code{busy} output raised when a word is taken and
dropped when the sum goes out, and its netlist agrees with its run
under nvc and Verilator, stalls included.

\begin{figure}[htbp]
\centering
\input{hold_timing.tex}
\caption{The stepper with held outputs: \code{phase} moves with the
state and stays at two through the wait for \code{go}; \code{busy}
is up from the word to the sum.}
\label{fig:x-hold}
\end{figure}

\lstinputlisting[caption={\code{ex\_hold.rs}. Run with \code{bazel run //lib/examples:ex\_hold}.},label={lst:x-hold}]{lib/examples/ex_hold.rs}

\lstinputlisting[style=ascii,caption={What \code{ex\_hold} printed when the build ran it: the stepper cycle by cycle, then the Verilog, with a register behind each output and the wire that reads ahead of it.},label={lst:x-hold-out}]{out_hold.txt}
